% =====================================================================
\part{Thermique architecturale}
% =====================================================================

\chapter{Régime dynamique}

\chapter{Enveloppe}

\chapter{Confort}
\label{Confort}

\section{Confort d'hiver}

\section{Confort d'été}
\subsection*{Les brasseurs d'air}
Un brasseur d'air peut fonctionner en mode été ou en mode hiver.
% [Figures : brasseur d'air en mode été / mode hiver]
Pour aller plus loin, voir les études de l'ADEME\footnote{ADEME, \emph{BRASSE : brasseurs d'air, une solution de sobriété et d'efficacité}, 2024 ; ADEME, \emph{BRASSE II : base de données des performances de 50 brasseurs d'air plafonniers à pales}, 2026. \texttt{https://librairie.ademe.fr}}.


\chapter{Équipements}

\section{Ventilation mécanique contrôlée (VMC)}

\section{Chauffage}

\section{Climatisation}

\chapter{Conclusion sur la thermique architecturale}
